\documentclass[10pt,a4paper]{article}
\usepackage[utf8]{inputenc}
\usepackage[T1]{fontenc}
\usepackage[spanish,es-nodecimaldot]{babel}
\usepackage{lmodern}
\usepackage{geometry}
\geometry{margin=1.8cm}
\usepackage{amsmath,amssymb,bm}
\usepackage{booktabs,longtable,array}
\usepackage{xcolor}
\usepackage{enumitem}
\usepackage{listings}
\usepackage[hidelinks]{hyperref}
\usepackage{microtype}
\usepackage{siunitx}

\definecolor{codegray}{gray}{0.96}
\lstset{
  basicstyle=\ttfamily\small,
  backgroundcolor=\color{codegray},
  frame=single,
  breaklines=true,
  columns=fullflexible,
  showstringspaces=false
}

\newcommand{\repo}{\texttt{PNNN\_PNGMP\_DOMP\_fair\_comparison}}
\newcommand{\code}[1]{\texttt{#1}}
\newcommand{\R}{\mathbb{R}}
\newcommand{\C}{\mathbb{C}}

\title{Especificación matemática y auditoría de implementación\\
del pipeline Complex GMP-DOMP--PN-IQ-GMP--PNNN}
\author{Documento interno previo a la redacción del paper}
\date{Commit \texttt{675687dc9b5a4b5c603a533f8b39ef5a53413eac}}

\begin{document}
\maketitle
\tableofcontents
\newpage

\section{Objetivo y alcance}

Este documento describe, en orden de ejecución y con correspondencia entre
código y matemáticas, el pipeline científico canónico del repositorio \repo.
Su finalidad es que los autores puedan justificar cada decisión durante la
redacción y la defensa del trabajo.

La auditoría cubre el código propio invocado por
\code{main\_sweep\_and\_comparison.m}, la generación de resultados, el
experimento widely-linear auxiliar y las pruebas directamente relacionadas.
El submódulo externo \code{third\_party/matlab2tikz} se trata como dependencia
y no se audita línea por línea. Los ficheros MAT binarios se consideran
artefactos de entrada o checkpoint; se auditan sus contratos de uso.

\section{Arquitectura del repositorio}

El punto de entrada es \code{main\_sweep\_and\_comparison.m}. La ejecución
canónica se organiza en:

\begin{enumerate}[leftmargin=*]
\item configuración central en \code{config/getFairDOMPComparisonConfig.m};
\item carga, preprocesado y split común;
\item barrido Complex GMP y PN-IQ-GMP;
\item referencias Fixed Ridge;
\item entrenamiento denso y pruning PNNN;
\item selección del punto operativo;
\item evaluación detallada y exportación de figuras.
\end{enumerate}

El código externo de matlab2tikz se incorpora mediante submódulo. Los MAT se
gestionan mediante Git LFS.

\section{Configuración canónica}

La configuración principal fija:

\begin{itemize}[leftmargin=*]
\item mapeo local \code{xy\_forward};
\item identificación: 10\% de la señal;
\item train interno: 85\% de identificación;
\item 10 bins internos de amplitud;
\item \(Q_p^{\max}=Q_n^{\max}=50\), \(P_{\max}=13\);
\item rejilla auxiliar \(\lambda=\{0,10^{-6},10^{-5},10^{-4},10^{-3},10^{-2}\}\), no usada por el sweep lineal canónico;
\item presupuestos \(P=20,30,\ldots,500\);
\item PNNN: \(M=13\), órdenes \(\{1,3,5,7\}\), modo full y \(H=12\);
\item punto de selección: ventana de 100 parámetros y tolerancia de
\(0.20\,\mathrm{dB}\).
\end{itemize}

Fuente: \code{config/getFairDOMPComparisonConfig.m}.

\section{Entrada, mapeo y eliminación de DC}

El fichero de medida aporta dos señales complejas \(x\) e \(y\). Con
\code{xy\_forward}:

\[
x_{\mathrm{in}}=x,\qquad y_{\mathrm{out}}=y.
\]

El comentario del código indica que son variables locales del bloque modelado;
el nombre no debe reinterpretarse automáticamente como dirección física del PA.

Antes de construir el split se elimina la media de toda la señal:

\[
x\leftarrow x-\frac{1}{N}\sum_n x[n],\qquad
y\leftarrow y-\frac{1}{N}\sum_n y[n].
\]

Fuentes:
\code{toolbox/pn\_gmp\_comparison/selectXYByMapping.m} y
\code{run\_parameter\_sweep.m}.

\section{Identificación, train interno, validation y señal completa}

\subsection{Selector de identificación}

La función \code{sel\_indices} localiza:

\[
n_{\max}=\arg\max_n |y[n]|
\]

y escoge un bloque contiguo de longitud aproximada \(0.1N\) alineado con
segmentos de ese tamaño y que contiene el máximo. La identificación no es una
muestra aleatoria de todo el registro.

\subsection{Split interno equilibrado en amplitud}

Sea \(\mathcal I\) la identificación. Se ordena \(\mathcal I\) por
\(|y[n]|\), se divide en diez bins y se asigna dentro de cada bin el 85\% a
train y el resto a validation usando una semilla fija. Este split se usa solo
para seleccionar entrenamiento y fine-tuning de la PNNN; los modelos lineales
usan toda \(\mathcal I\).

\subsection{Dominio de evaluación}

\[
\mathcal F=\{1,\ldots,N\}.
\]

Como \(\mathcal I\subset\mathcal F\), la señal completa no es un holdout
independiente. Debe llamarse full-signal evaluation, no independent test.

Fuente: \code{toolbox/splits/buildCommonComparisonSplit.m}.

\section{Población GMP}

\subsection{Representación genérica del objeto Regressor}

Cada columna se representa mediante tres multisets de retardos:

\[
\mathcal X_k,\qquad \mathcal X_k^*,\qquad \mathcal E_k.
\]

La columna es:

\[
u_k[n]
=
\prod_{q\in\mathcal X_k}x[n-q]
\prod_{q\in\mathcal X_k^*}x^*[n-q]
\prod_{q\in\mathcal E_k}|x[n-q]|.
\]

Las repeticiones de un mismo retardo implementan potencias. El método
\code{deriveEnvelopeTerms} convierte pares \(x[n-q]x^*[n-q]\) en
\(|x[n-q]|^2\).

\subsection{Familias GMP inicializadas}

Con \(P_{\mathrm{GMP}}=13\), \(L=10\) y \(M=2\):

\[
u^{(a)}_{k,\ell}[n]
=x[n-\ell]|x[n-\ell]|^k,
\quad
k=0,\ldots,12,\ \ell=0,\ldots,10,
\]

\[
u^{(b)}_{k,\ell,m}[n]
=x[n-\ell]|x[n-\ell-m]|^k,
\quad
k=1,\ldots,12,\ m=1,2,
\]

\[
u^{(c)}_{k,\ell,m}[n]
=x[n-\ell]|x[n-\ell+m]|^k.
\]

El código incluye todos los grados totales del 1 al 13. Después se añaden
provisionalmente \(x[n]\), \(x^*[n]\) y \(|x[n]|\) y se eliminan duplicados.
El tamaño esperado es 673 columnas: 671 GMP canónicas y dos auxiliares no
duplicados.

\subsection{Indexado}

La construcción usada por el sweep aplica extensión periódica:

\[
x[n-q]=x[((n-q-1)\bmod N)+1].
\]

Fuente:
\code{toolbox/pn\_gmp\_comparison/regressorManager.m},
\code{Regressor.m},
\code{modelsConfig.m} y
\code{buildGMPRegressorRows.m}.

\section{Selección DOMP}

\subsection{Entrada y salida}

\code{selectDOMPSupport(X,y,K,\tau)} recibe una matriz candidata \(X\), un
objetivo real o complejo \(y\), el máximo de componentes y una tolerancia de
columna.

\subsection{Recurrencia}

Inicialmente:

\[
Z_0=X,\qquad r_0=y,\qquad\mathcal S_0=\varnothing.
\]

En la iteración \(t\), para las columnas elegibles:

\[
\bar z_j=\frac{z_j}{\|z_j\|_2},
\qquad
s_j=\|\bar z_j^H r_{t-1}\|_2.
\]

Se selecciona:

\[
j_t=\arg\max_j s_j.
\]

Con \(q=\bar z_{j_t}\), se actualiza:

\[
Z_t=Z_{t-1}-q(q^H Z_{t-1}).
\]

Después se ignora la matriz ortogonalizada para el ajuste y se resuelve sobre
las columnas originales:

\[
\hat c_t
=\arg\min_c\|y-X_{\mathcal S_t}c\|_2^2,
\qquad
r_t=y-X_{\mathcal S_t}\hat c_t.
\]

El solver usa QR cuando el rango se considera completo y
\code{lsqminnorm} en caso contrario.

\subsection{Aclaraciones esenciales}

\begin{itemize}[leftmargin=*]
\item DOMP selecciona \emph{columnas}, no ``coeficientes''.
\item Ridge no participa en la selección del soporte.
\item Complex GMP y PN-IQ tienen rutas DOMP distintas.
\item Cada ruta máxima se calcula una vez y los presupuestos usan prefijos.
\item En PN-IQ, \(X\) es real y \(y\) complejo; el módulo de la correlación
combina ambas componentes de salida.
\end{itemize}

Fuente: \code{toolbox/domp/selectDOMPSupport.m}.

\section{Complex GMP-DOMP}

\subsection{Ruta de identificación}

Se construye \(U\) sobre toda la identificación. DOMP obtiene una única ruta
máxima de 250 regresores complejos. Para \(P\) parámetros reales:

\[
K=P/2,\qquad
\mathcal S_P=\{s_1,\ldots,s_K\}.
\]

\subsection{Normalización y lambda}

Sea:

\[
D=\mathrm{diag}(\max_n|u_1[n]|,\ldots,\max_n|u_K[n]|),
\qquad
\widetilde U=UD^{-1}.
\]

Con \(p_y=\max_n|y[n]|\), se ajusta la salida de pico unitario
\(\widetilde y=y/p_y\).

Para \(\lambda=0\):

\[
\widetilde h_0=\widetilde U^\dagger \widetilde y.
\]

Para \(\lambda>0\):

\[
\widetilde h_\lambda
=(\widetilde U^H\widetilde U+\lambda I)^{-1}
\widetilde U^H\widetilde y.
\]

Los coeficientes en la matriz física son:

\[
h_\lambda=p_yD^{-1}\widetilde h_\lambda.
\]

La solución principal usa siempre \(\lambda=0\) y se calcula mediante
\code{lsqminnorm} con tolerancia de rango robusta. La rejilla
\code{cfg.lambdaGrid} queda reservada a experimentos auxiliares.

\subsection{Evaluación}

La única ruta se selecciona sobre toda identificación. Todos sus prefijos se
ajustan con \(\lambda=0\) y se evalúan en identificación y señal completa.

\subsection{Parámetros}

\[
P_{\mathrm{GMP}}=2K.
\]

Fuente: \code{toolbox/sweep/fit\_complex\_gmp\_domp.m}.

\section{Modelo PN-IQ-GMP propuesto}

\subsection{Normalización de fase}

\[
r[n]=
\begin{cases}
\dfrac{x^*[n]}{|x[n]|},&|x[n]|\neq0,\\[4pt]
1,&|x[n]|=0.
\end{cases}
\]

El objetivo y cada regresor se expresan en el marco de fase de la entrada:

\[
z[n]=r[n]y[n],
\qquad
\widetilde u_j[n]=r[n]u_j[n].
\]

\subsection{Diccionario I/Q}

\[
f_{j,I}[n]=\Re\{\widetilde u_j[n]\},
\qquad
f_{j,Q}[n]=\Im\{\widetilde u_j[n]\}.
\]

Para términos canónicos con portador \(x[n]\), \(f_{j,Q}=0\) y la columna se
elimina. Los auxiliares se rotan directamente. El resultado es una matriz real
\(F\).

\subsection{Selección}

DOMP se aplica a \(F_{\mathrm{tr}}\in\R^{N_{\mathrm{tr}}\times J}\) y al
objetivo complejo \(z_{\mathrm{tr}}\in\C^{N_{\mathrm{tr}}}\). Para un
presupuesto \(P\):

\[
M=P/2.
\]

El soporte contiene \(M\) características reales, no \(M\) regresores
complejos.

\subsection{Ajuste independiente}

Sobre el mismo soporte:

\[
\widetilde c_I
=\arg\min_c
\|\Re\{z\}-\widetilde F_{\mathcal S}c\|_2^2
+\lambda\|c\|_2^2,
\]

\[
\widetilde c_Q
=\arg\min_c
\|\Im\{z\}-\widetilde F_{\mathcal S}c\|_2^2
+\lambda\|c\|_2^2.
\]

La predicción rotada y la restauración son:

\[
\hat z=F_{\mathcal S}c_I+jF_{\mathcal S}c_Q,
\qquad
\hat y[n]=r^*[n]\hat z[n].
\]

El soporte es común a los dos ajustes. Ambos se resuelven independientemente
mediante \code{lsqminnorm} y usan \(\lambda=0\) en la solución principal.
Para el ajuste, cada característica activa se divide por su máximo absoluto y
la salida rotada se divide por el pico de la salida original. La transformación
inversa se aplica únicamente a los coeficientes usados para predecir.

\subsection{Libertad adicional}

Para una columna compleja:

\[
\widetilde u=a+jb.
\]

Un coeficiente complejo convencional impone:

\[
h\widetilde u=h\,a+(jh)\,b.
\]

Es decir, si se interpreta mediante dos características reales, sus
coeficientes complejos están ligados:

\[
g_b=jg_a.
\]

PN-IQ permite:

\[
\hat z=g_a a+g_b b,
\qquad
g_a,g_b\in\C,
\]

sin esa restricción, y puede seleccionar \(a\) y \(b\) por separado.

\subsection{Contabilidad}

Cada característica seleccionada tiene un coeficiente para I y otro para Q:

\[
P_{\mathrm{PN-IQ}}=2M.
\]

El nombre recomendado es \textbf{PN-IQ-GMP}. DOMP es el método estándar
de selección aplicado al diccionario normalizado en fase.

Fuente: \code{toolbox/sweep/fit\_independent\_pniq\_domp.m} y
\code{factorizeGMPRegressor.m}.

\section{Fixed Ridge}

Las variantes fijas conservan las rutas DOMP principales y repiten el ajuste
con:

\[
\lambda\in\{10^{-3},10^{-4},10^{-5}\}.
\]

No vuelven a seleccionar soporte. Por contrato, parámetros y FLOPs son los
mismos que en la fila principal correspondiente. Solo cambian coeficientes,
NMSE y rango numérico.

Fuente: \code{toolbox/sweep/run\_fixed\_ridge\_sweep.m}.

\section{Dataset PNNN}

\subsection{Características}

Con \(M=13\), se toman:

\[
x_n=[x[n],x[n-1],\ldots,x[n-13]]
\]

con extensión periódica. El vector full es:

\[
\phi[n]=
\begin{bmatrix}
\Re\{r[n]x_n\}\\
\Im\{r[n]x_n\}\\
|x_n|^1\\
|x_n|^3\\
|x_n|^5\\
|x_n|^7
\end{bmatrix}
\in\R^{84}.
\]

El objetivo es:

\[
t[n]=
\begin{bmatrix}
\Re\{r[n]y[n]\}\\
\Im\{r[n]y[n]\}
\end{bmatrix}.
\]

\subsection{Redundancias explícitas}

\[
\Im\{r[n]x[n]\}=0,
\qquad
\Re\{r[n]x[n]\}=|x[n]|.
\]

Por ello, el modo full contiene una entrada estructuralmente cero y una
duplicación con el término de envolvente de orden uno del tap actual.

Fuente:
\code{toolbox/pnnn/buildPhaseNormDataset.m}.

\section{Arquitectura y entrenamiento PNNN}

\subsection{Normalización}

\[
\bar\phi_d=\frac{\phi_d-\mu_{X,d}}{\sigma_{X,d}},
\qquad
\bar t_q=\frac{t_q-\mu_{Y,q}}{\sigma_{Y,q}}.
\]

Los \(\sigma=0\) se sustituyen por uno.

\subsection{Red}

\[
h=\sigma(W_1\bar\phi+b_1),
\qquad
\bar{\hat t}=W_2h+b_2,
\]

con \(D=84\), \(H=12\) y salida 2. La activación oculta es sigmoide y la
salida es lineal.

\[
P_{\mathrm{dense}}
=DH+H+2H+2
=H(D+3)+2
=1046.
\]

\subsection{Selección de épocas}

La red densa de selección se entrena con train interno y validation. Se
conserva la red de menor validation loss y se convierte la iteración óptima en
épocas.

El número máximo de épocas, el periodo de caída y la paciencia se escalan para
conservar aproximadamente el número histórico de actualizaciones definido por
70\% de la señal y 150 épocas.

\subsection{Refit}

Tras seleccionar la época, se reinicia la red con la misma semilla y se
entrena desde cero sobre toda identificación. Las estadísticas z-score se
recalculan en identificación.

Fuente:
\code{prepare\_pnnn\_dense\_source.m}.

\section{Pruning PNNN}

\subsection{Máscara}

Se ordenan globalmente por magnitud todos los pesos de \(W_1\) y \(W_2\). Los
sesgos quedan siempre activos. Si el objetivo es \(P\):

\[
P_{\mathrm{bias}}=12+2=14,
\qquad
P_{\mathrm{weights}}=P-14.
\]

\subsection{Selección de fine-tuning}

La duración de fine-tuning se selecciona una sola vez podando la red interna
al presupuesto 200. Esa duración se transfiere a todos los puntos finales.

\subsection{Refit sparse}

Cada presupuesto:

\begin{enumerate}[leftmargin=*]
\item parte de la misma red densa final inmutable;
\item construye su máscara por magnitud;
\item aplica máscara a parámetros;
\item aplica máscara a gradientes;
\item actualiza con Adam;
\item reaplica la máscara tras cada actualización;
\item verifica que los pesos podados siguen en cero;
\item fine-tunea durante la duración común sobre toda identificación.
\end{enumerate}

Las máscaras, antes de fine-tuning, son prefijos por magnitud de una misma
fuente; las redes finales se optimizan de forma independiente.

Fuentes:
\code{pruneAndFineTunePNNN.m} y
\code{fit\_sparse\_pnnn\_target.m}.

\section{Predicción PNNN}

La salida normalizada se desnormaliza:

\[
\hat t=\bar{\hat t}\odot\sigma_Y+\mu_Y.
\]

Se reconstruye:

\[
\hat z=\hat t_I+j\hat t_Q,
\qquad
\hat y[n]=r^*[n]\hat z[n].
\]

Fuente: \code{toolbox/pnnn/predictPhaseNorm.m}.

\section{Métrica NMSE}

Todos los modelos usan:

\[
\mathrm{NMSE}_{\mathrm{dB}}
=10\log_{10}
\left(
\frac{\sum_n|y[n]-\hat y[n]|^2}
{\sum_n|y[n]|^2}
\right).
\]

Fuente: \code{toolbox/metrics/nmseComplexDb.m}.

\section{Complejidad}

\subsection{Convención}

\begin{center}
\begin{tabular}{lc}
\toprule
Operación & FLOPs\\
\midrule
Suma real & 1\\
Multiplicación real & 1\\
MAC real & 2\\
Suma compleja & 2\\
Multiplicación compleja & 6\\
MAC complejo & 8\\
Suma de bias real & 1\\
\bottomrule
\end{tabular}
\end{center}

Raíces, divisiones, valores absolutos y activaciones se reportan como
operaciones especiales sin peso FLOP.

\subsection{Lineales}

\code{countModelOperations} separa generación de características,
multiplicaciones por coeficientes, acumulación, normalización y restauración de
fase. PN-IQ usa un schedule structure-aware que reutiliza magnitudes y
portadores rotados.

\subsection{PNNN}

El coste denso incluye construcción de 84 características, dos capas, biases,
12 sigmoides y restauración. El coste sparse sustituye el número de pesos
densos por pesos activos, pero mantiene todo el coste de generación de
características.

El valor principal sparse presupone que una implementación omite exactamente
los pesos cero. El \code{predict} MATLAB actual no constituye esa
implementación sparse.

Fuentes:
\code{countModelOperations.m},
\code{countModelFLOPs.m},
\code{getFLOPConvention.m},
\code{countPNNNFLOPs.m} y
\code{countSparsePNNNFLOPs.m}.

\section{Rango de coeficientes}

\subsection{Construcción lineal equivalente}

Cada columna se normaliza por su pico:

\[
\bar u_k=\frac{u_k(x)}{\max_n|u_k(x)[n]|}.
\]

La salida de identificación se normaliza a pico unitario y se ajusta en esa
base. Para GMP:

\[
C_{\max}^{\mathrm{GMP}}
=\max_k|\bar h_k|.
\]

Para PN-IQ:

\[
C_{\max}^{\mathrm{PN-IQ}}
=\max_k\big(
|\bar c_{I,k}|,
|\bar c_{Q,k}|
\big).
\]

Una escala global positiva de entrada es equivalente, porque su factor cancela
al dividir cada columna por su propio pico. La solución en la base de pico
unitario define la métrica y se transforma únicamente para predecir la salida
original. Las pruebas lineales reconstruyen la salida y las columnas de pico
unitario, y verifican igualdad.

\subsection{PNNN}

La PNNN reporta:

\[
C_{\max}^{\mathrm{PNNN}}
=\max_{\theta_i\ \mathrm{activo}}|\theta_i|
\]

en la parametrización z-score, incluyendo biases. No es la misma
parametrización que la de los modelos lineales.

\section{Selección del punto operativo}

Para cada presupuesto \(P\) y familia \(m\):

\[
G_m(P)=
\mathrm{NMSE}_m(P)
-
\min_{P\le q\le P+100}
\mathrm{NMSE}_m(q).
\]

Se selecciona el menor \(P\) con ventana completa para el que:

\[
G_m(P)\le0.20\,\mathrm{dB}
\quad\forall m.
\]

El resultado es \(P=340\).

Fuente: \code{toolbox/sweep/selectOperatingPoint.m}.

\section{Espectros y figuras temporales}

\subsection{Welch}

\[
e_m[n]=y[n]-\hat y_m[n].
\]

Todas las señales y errores se procesan conjuntamente con:

\begin{itemize}[leftmargin=*]
\item Hann periódica, 2048 muestras;
\item 75\% de solapamiento;
\item \(N_{\mathrm{FFT}}=16384\);
\item espectro bilateral centrado;
\item referencia: máximo de PSD de la salida medida;
\item suelo numérico: \(-150\,\mathrm{dB}\).
\end{itemize}

\subsection{Tiempo}

Se seleccionan las 1000 muestras centrales. Se exportan dos figuras:

\begin{enumerate}[leftmargin=*]
\item magnitud + parte real;
\item magnitud + parte imaginaria.
\end{enumerate}

La superposición visual significa que los residuos son pequeños frente a la
señal, no que los NMSE sean iguales.

Fuentes:
\code{computeSelectedPointSpectra.m} y
\code{run\_selected\_comparison.m}.

\section{Widely-linear GMP auxiliar}

El probe construye:

\[
U_{\mathrm{WL}}=[U,\ U^*]
\]

y selecciona 170 átomos con DOMP para conservar 340 parámetros reales. No
forma parte del sweep canónico y fija actualmente el directorio de resultados.

\begin{center}
\begin{tabular}{lrr}
\toprule
Modelo & Validation NMSE (dB) & Mejora vs. GMP\\
\midrule
Complex GMP-DOMP & -38.85998 & 0\\
WL-GMP & -38.86672 & 0.00674\\
PN-IQ-GMP & -42.74017 & 3.88019\\
\bottomrule
\end{tabular}
\end{center}

El resultado muestra que concatenar \(U^*\) no reproduce la ganancia de
PN-IQ.

\section{Resultados auditados en 340 parámetros}

\begin{center}
\begin{tabular}{lrrr}
\toprule
Familia & Validation NMSE (dB) & FLOPs/muestra & Máx. escalar\\
\midrule
Complex GMP-DOMP & -38.85998 & 1807 & \(6.34896\times10^6\)\\
PN-IQ-GMP & -42.74017 & 1095 & \(4.59256\times10^6\)\\
PNNN & -39.24296 & 840 & 2.17243\\
\bottomrule
\end{tabular}
\end{center}

PN-IQ mejora \(3.88019\,\mathrm{dB}\) respecto a GMP. Las referencias Ridge
reducen el rango numérico lineal a aproximadamente \(86\)--\(132\), con
pérdida de NMSE.

\section{Checkpoints y reproducibilidad}

La identidad SHA-256 del sweep incorpora:

\begin{itemize}[leftmargin=*]
\item firma de experimento;
\item rejilla de parámetros, lambdas Ridge fijas y protocolo lineal;
\item especificación de población GMP;
\item tolerancia DOMP;
\item configuración PNNN;
\item entrenamiento y pruning.
\end{itemize}

Los checkpoints lineal, Ridge, denso y sparse se rechazan si la identidad no
coincide. La red densa se firma además con sus learnables y estadísticas de
normalización.

\section{Hallazgos y acciones}

\subsection{Crítico: \code{clear} dentro del main}

El main actual comienza con:

\begin{lstlisting}[language=Matlab]
clear; clc; close all force;
\end{lstlisting}

\code{clear} elimina las variables del workspace de la función y puede borrar
el argumento \code{selectedParameters}. Debe quedar, como máximo:

\begin{lstlisting}[language=Matlab]
clc;
close all force;
\end{lstlisting}

o no ejecutar ninguna limpieza dentro de la función.

\subsection{Nombre del método}

Sustituir:

La nomenclatura canónica del modelo normalizado en fase es
\textbf{PN-IQ-GMP}; DOMP describe exclusivamente la selección del soporte.

por:

\begin{itemize}[leftmargin=*]
\item modelo: \textbf{PN-IQ-GMP};
\item identificación: \textbf{DOMP-based sparse support selection};
\item leyenda compacta: \textbf{PN-IQ-GMP}.
\end{itemize}

\subsection{Limitaciones que deben declararse}

\begin{enumerate}[leftmargin=*]
\item full-signal incluye identificación;
\item PNNN y lineales no comparten parametrización para \(C_{\max}\);
\item PNNN sparse FLOPs presupone kernel sparse;
\item los espacios candidatos difieren;
\item la PNNN full contiene una entrada cero y una duplicada;
\item WL probe tiene path de resultados hardcoded.
\end{enumerate}

\section{Preguntas de defensa y respuestas}

\subsection*{¿Sobre qué se aplica DOMP?}
Sobre columnas candidatas. En GMP, columnas complejas. En PN-IQ, columnas
reales I/Q. Nunca se aplica sobre los coeficientes ya ajustados.

\subsection*{¿Cuándo se ejecuta DOMP?}
Una sola vez por familia lineal sobre toda identificación. Cada presupuesto
usa un prefijo de esa ruta.

\subsection*{¿Comparten soporte I y Q?}
Sí. Las dos salidas usan el mismo conjunto de características reales; sus
coeficientes se estiman por separado.

\subsection*{¿La PNNN usa DOMP?}
No. Se entrena densa y se poda globalmente por magnitud.

\subsection*{¿Las configuraciones podadas de PNNN forman una trayectoria iterativa?}
No. Todos parten de una misma fuente densa final y se fine-tunean
independientemente.

\subsection*{¿Por qué PN-IQ es más flexible?}
Porque elimina la restricción \(g_b=jg_a\) entre los pesos de las componentes
real e imaginaria de un regressor rotado y permite seleccionarlas por separado.

\subsection*{¿Es independiente el NMSE full-signal?}
No. Incluye identificación.

\subsection*{¿Los FLOPs son tiempo de ejecución?}
No. Son cuentas analíticas de sumas y multiplicaciones bajo una convención
fija; las operaciones especiales quedan separadas.

\appendix
\section{Mapa código--concepto}

\begin{longtable}{p{0.38\textwidth}p{0.55\textwidth}}
\toprule
Concepto & Implementación\\
\midrule
Configuración & \code{config/getFairDOMPComparisonConfig.m}\\
Entrada canónica & \code{main\_sweep\_and\_comparison.m}\\
Pipeline & \code{run\_parameter\_sweep.m}\\
Split & \code{toolbox/splits/buildCommonComparisonSplit.m}\\
Selector 10\% & \code{toolbox/pn\_gmp\_comparison/sel\_indices.m}\\
Población GMP & \code{GMP\_createRegressorManager.m}, \code{regressorManager.m}\\
Evaluación GMP & \code{buildGMPRegressorRows.m}\\
DOMP & \code{toolbox/domp/selectDOMPSupport.m}\\
Complex GMP-DOMP & \code{toolbox/sweep/fit\_complex\_gmp\_domp.m}\\
PN-IQ-GMP & \code{toolbox/sweep/fit\_pniq\_gmp.m}\\
Ridge & \code{toolbox/sweep/run\_fixed\_ridge\_sweep.m}\\
Features PNNN & \code{buildPhaseNormDataset.m}\\
Fuente densa y normalización & \code{prepare\_pnnn\_dense\_source.m}\\
Máscaras y fine-tuning & \code{pruneAndFineTunePNNN.m}\\
Sparse refit & \code{fit\_sparse\_pnnn\_target.m}\\
Predicción & \code{predictPhaseNorm.m}\\
NMSE & \code{nmseComplexDb.m}\\
Complejidad & \code{countModelOperations.m}, \code{countModelFLOPs.m}\\
Selección 340 & \code{selectOperatingPoint.m}\\
Espectro & \code{computeSelectedPointSpectra.m}\\
Exportación & \code{exportPaperFigure.m}\\
Probe WL & \code{run\_widely\_linear\_gmp\_probe.m}\\
\bottomrule
\end{longtable}

\section{Glosario de nombres recomendado}

\begin{description}[leftmargin=4cm,style=nextline]
\item[Complex GMP-DOMP] Baseline complejo con soporte seleccionado mediante DOMP.
\item[PN-IQ-GMP] Modelo propuesto.
\item[DOMP-selected PN-IQ-GMP] Forma descriptiva completa.
\item[PNNN] Benchmark neural de 12 neuronas ocultas.
\item[Validation NMSE (dB)] NMSE sobre todo el registro, incluyendo identificación.
\item[Identification NMSE] NMSE sobre el bloque del 10\%.
\item[Internal validation] Subconjunto de identificación usado solo por la PNNN.
\end{description}

\end{document}
